\begin{proof}[Proof of \cref{obj:lem:causal-identification}]
\begin{enumerate}
\item Let \(P_c\) be the joint causal law on
\(\mathbb R^d\times\{0,1\}\times\mathbb R^3\):
\[
P_c:=\mathcal L\bigl(X,A,Y(0),Y(1),Y\bigr).
\]
Its observed marginal is \(P\), and its covariate marginal is \(P_X\). Define the joint conditioning law and conditional probability kernels by
\[
\begin{aligned}
\pi_{X,A}&:=\mathcal L_{P_c}(X,A),\\
K^c_{A\mid X}(x)&:=\mathcal L_{P_c}(A\mid X=x),\\
K_{1\mid X}(x)&:=\mathcal L_{P_c}(Y(1)\mid X=x),\\
K^c_{Y\mid X,A}(x,a)&:=\mathcal L_{P_c}(Y\mid X=x,A=a),\\
K^c_{1\mid X,A}(x,a)&:=\mathcal L_{P_c}(Y(1)\mid X=x,A=a),
\end{aligned}
\qquad x\in\mathbb R^d,\quad a\in\{0,1\}.
\]
Choose measurable versions defined at every conditioning input, with arbitrary probability values on exceptional conditioning sets.

By \cref{obj:ass:consistency},
\[
A=1\ \Longrightarrow\ Y=Y(1)
\qquad P_c\text{-almost surely}.
\]
Conditioning this equality on \((X,A)\) gives
\[
a=1\ \Longrightarrow\
K^c_{Y\mid X,A}(x,a)=K^c_{1\mid X,A}(x,a)
\qquad \pi_{X,A}\text{-almost everywhere}.
\]
By \cref{obj:ass:exchangeability}, \(Y(1)\) and \(A\) are conditionally independent given \(X\). Consequently,
\[
K^c_{1\mid X,A}(x,a)=K_{1\mid X}(x)
\qquad \pi_{X,A}\text{-almost everywhere}.
\]
% lean: CausalSmith.Stat.WeakOverlap.mu1_eq_treatedRegression (hfiber, hci, hexch)

\item Applying \cref{obj:ass:global-tail} at \(t=0\), and using \(\gamma-1>0\), yields
\[
0\le P_X\{e\le0\}\le C\,0^{\gamma-1}=0.
\]
Since \(\{e=0\}\subseteq\{e\le0\}\),
\[
P_X\{e=0\}=0.
\]
To express positivity through the conditioning kernel, define
\[
K^o_{A\mid X}(x):=\mathcal L_P(A\mid X=x),
\qquad x\in\mathbb R^d,
\]
again as a measurable probability kernel defined everywhere. The observed marginal relation and the definition of the propensity give
\[
K^o_{A\mid X}(x)=K^c_{A\mid X}(x),
\qquad
e(x)=K^o_{A\mid X}(x)(\{1\})\ge0
\qquad P_X\text{-almost everywhere}.
\]
Combining these identities with the null-set calculation proves
\[
K^c_{A\mid X}(x)(\{1\})=e(x)>0
\qquad P_X\text{-almost everywhere}.
\]
% lean: CausalSmith.Stat.WeakOverlap.positivePropensity_ae

\item Disintegration of the joint conditioning law gives
\[
\pi_{X,A}(dx,da)=P_X(dx)\,K^c_{A\mid X}(x,da).
\]
We use the following consequence of this identity. For any predicate
\(p:\mathbb R^d\times\{0,1\}\to\{\text{true},\text{false}\}\) holding
\(\pi_{X,A}\)-almost everywhere, disintegration implies
\[
K^c_{A\mid X}(x)\{a:\neg p(x,a)\}=0
\qquad P_X\text{-almost everywhere}.
\]
At any such \(x\) with \(K^c_{A\mid X}(x)(\{1\})>0\), failure of \(p(x,1)\) would imply
\[
0<K^c_{A\mid X}(x)(\{1\})
\le K^c_{A\mid X}(x)\{a:\neg p(x,a)\}=0.
\]
Thus \(p(x,1)\) holds \(P_X\)-almost everywhere. Applying this argument to the conjunction of the two conditional-kernel identities above gives
\[
K^c_{Y\mid X,A}(x,1)=K_{1\mid X}(x)
\qquad P_X\text{-almost everywhere}.
\]

Now define the observed outcome kernel and its treated slice by
\[
\begin{aligned}
K^o_{Y\mid X,A}(x,a)&:=\mathcal L_P(Y\mid X=x,A=a),\\
K_{\mathrm{tr}}(x)&:=K^o_{Y\mid X,A}(x,1),
\end{aligned}
\qquad x\in\mathbb R^d,\quad a\in\{0,1\}.
\]
The observed marginal relation gives
\[
K^o_{Y\mid X,A}(x,a)=K^c_{Y\mid X,A}(x,a)
\qquad \pi_{X,A}\text{-almost everywhere}.
\]
Applying the same positive-slice argument once more, we obtain
\[
K_{\mathrm{tr}}(x)
=K^c_{Y\mid X,A}(x,1)
=K_{1\mid X}(x)
\qquad P_X\text{-almost everywhere}.
\]
% lean: CausalSmith.Stat.WeakOverlap.ae_slice_of_ae_compProd_of_pos

\item Define the treated regression on all of \(\mathbb R^d\) by
\[
r_{\mathrm{tr}}(x):=\int_{\mathbb R}y\,K_{\mathrm{tr}}(x,dy),
\]
using value zero wherever the integrand is not integrable. By
\cref{obj:ass:bounded-potential-outcomes} and the kernel equality just proved, for \(P_X\)-almost every \(x\),
\[
|r_{\mathrm{tr}}(x)|\le B,
\qquad
\int_{\mathbb R}
e^{\lambda(y-r_{\mathrm{tr}}(x))}\,K_{1\mid X}(x,dy)
\le e^{B^2\lambda^2/2}
\quad\text{for every }\lambda\in\mathbb R,
\]
and each exponential integrand in this display is integrable.

The choices \(\lambda=1\) and \(\lambda=-1\) control the absolute first moment. Indeed, for every \(y\in\mathbb R\),
\[
1+|y-r_{\mathrm{tr}}(x)|
\le e^{|y-r_{\mathrm{tr}}(x)|}
\le e^{y-r_{\mathrm{tr}}(x)}
   +e^{-(y-r_{\mathrm{tr}}(x))},
\]
so
\[
|y|
\le |y-r_{\mathrm{tr}}(x)|+|r_{\mathrm{tr}}(x)|
\le e^{y-r_{\mathrm{tr}}(x)}
   +e^{-(y-r_{\mathrm{tr}}(x))}
   +|r_{\mathrm{tr}}(x)|.
\]
The right-hand side is integrable under \(K_{1\mid X}(x)\). Integrating the inequality and using the exponential bounds yields
\[
\begin{aligned}
\int_{\mathbb R}|y|\,K_{1\mid X}(x,dy)
&\le
\int_{\mathbb R}e^{y-r_{\mathrm{tr}}(x)}\,K_{1\mid X}(x,dy)
+\int_{\mathbb R}e^{-(y-r_{\mathrm{tr}}(x))}\,K_{1\mid X}(x,dy)
+|r_{\mathrm{tr}}(x)|\\
&\le 2e^{B^2/2}+B.
\end{aligned}
\]
This bound is uniform on a set of full \(P_X\)-measure. Disintegration therefore gives
\[
\int |Y(1)|\,dP_c
=
\int_{\mathbb R^d}
\left(\int_{\mathbb R}|y|\,K_{1\mid X}(x,dy)\right)P_X(dx)
\le 2e^{B^2/2}+B<\infty.
\]
In particular, \(Y(1)\) is integrable under \(P_c\).
% lean: CausalSmith.Stat.WeakOverlap.mu1_eq_treatedRegression (houtcome, hY₁fiber, hY₁norm_bound, hY₁int)

\item Define
\[
k_1(x):=\int_{\mathbb R}y\,K_{1\mid X}(x,dy),
\qquad x\in\mathbb R^d,
\]
with value zero wherever the integrand is not integrable. The conditional-distribution formula for conditional expectation, together with the defining conditional-mean relation for the treated causal response, gives
\[
\mu_1(X)
=\mathbb E_{P_c}[Y(1)\mid\sigma(X)]
=k_1(X)
\qquad P_c\text{-almost surely}.
\]
The kernel integral \(k_1\) is measurable. Continuity of \(\mu_1\) on \(\mathcal X\), supplied by
\cref{obj:ass:holder-response,obj:def:holder}, and
\(P_X\ll\operatorname{Leb}|_{\mathcal X}\) supply a measurable representative of \(\mu_1\) under \(P_X\). Using these representatives, the covariate pushforward identity transfers the preceding equality to
\[
\mu_1(x)=k_1(x)
\qquad P_X\text{-almost everywhere}.
\]
The equality \(K_{\mathrm{tr}}=K_{1\mid X}\) then gives
\[
\mu_1(x)
=k_1(x)
=r_{\mathrm{tr}}(x)
=\mathbb E_P[Y\mid A=1,X=x]
\qquad P_X\text{-almost everywhere}.
\]
% lean: CausalSmith.Stat.WeakOverlap.mu1_eq_treatedRegression (hce, hcomp, hbase, hident)

\item It remains to establish the asserted equality for every point of the cube. Define the restricted Lebesgue measure by
\[
\nu:=\operatorname{Leb}|_{\mathcal X}.
\]
The parameter restrictions give \(c_f>0\). Absolute continuity of \(P_X\) with respect to \(\nu\) and \cref{obj:ass:density} yield
\[
P_X=f\nu,
\qquad
f=\frac{dP_X}{d\nu}\ge c_f>0
\qquad \nu\text{-almost everywhere}.
\]
Consequently,
\[
\nu\{f=0\}=0,
\qquad
\nu\ll f\nu=P_X.
\]
For completeness, if a Borel set \(E\subseteq\mathbb R^d\) satisfies \(P_X(E)=0\), then
\[
\int_E f\,d\nu=P_X(E)=0.
\]
Since \(f\) is nonnegative, this forces \(f=0\) \(\nu\)-almost everywhere on \(E\); its strict positivity \(\nu\)-almost everywhere gives \(\nu(E)=0\).

Let \(g:\mathbb R^d\to\mathbb R\) satisfy the continuity and almost-everywhere agreement required in the statement. The identifying equality already established implies
\[
g(x)=\mu_1(x)
\qquad P_X\text{-almost everywhere},
\]
and hence, by \(\nu\ll P_X\),
\[
g(x)=\mu_1(x)
\qquad \nu\text{-almost everywhere}.
\]
Both functions are continuous on \(\mathcal X\). Thus
\[
g(x)=\mu_1(x)
\qquad\text{for every }x\in\operatorname{int}(\mathcal X):
\]
a nonzero difference at an interior point would, by continuity, persist on an open ball contained in the cube, whose positive Lebesgue measure contradicts the preceding almost-everywhere equality. Finally,
\[
\overline{\operatorname{int}(\mathcal X)}
=\prod_{i=1}^{d}\overline{(0,1)}
=\prod_{i=1}^{d}[0,1]
=\mathcal X.
\]
Continuity therefore extends the equality to every \(x\in\mathcal X\), as required.
% lean: CausalSmith.Stat.WeakOverlap.mu1_eq_treatedRegression (hreverse, hAEvol, hcube_reg)
\end{enumerate}
\end{proof}
